\documentclass[12pt,a4paper]{article}
\usepackage[margin=1in]{geometry}
\usepackage{amsmath,amssymb}
\usepackage{booktabs}
\usepackage{xcolor}
\usepackage{listings}
\usepackage{float}

\definecolor{codegray}{gray}{0.95}
\lstset{
backgroundcolor=\color{codegray},
basicstyle=\ttfamily\small,
frame=single,
breaklines=true,
columns=fullflexible
}

\title{\textbf{Clustering Human Activity Recognition Data}\\
using K-Means, DBSCAN, and Hierarchical Clustering}
\author{Nathaniel Christian\\Machine Learning\\Reg. No: 3122247001037}
\date{}

\begin{document}
\maketitle

\section{Aim}
To implement and analyze the performance of clustering algorithms on the Human Activity Recognition dataset:
\begin{itemize}
\item Model A: K-Means Clustering.
\item Model B: DBSCAN (Density-Based Spatial Clustering of Applications with Noise).
\item Model C: Hierarchical Agglomerative Clustering (HAC).
\end{itemize}

\section{Preprocessing Steps}
\begin{itemize}
\item Load the datasets.
\item Handle missing values.
\item Encode categorical labels.
\item Apply normalization or standardization.
\item Perform exploratory data analysis (EDA).
\item Visualize feature distributions.
\item Apply dimensionality reduction using PCA/t-SNE for visualization.
\end{itemize}

\section{Code}
\begin{lstlisting}[language=Python]
import numpy as np
import pandas as pd
import matplotlib.pyplot as plt
from sklearn.preprocessing import StandardScaler
from sklearn.decomposition import PCA
from sklearn.cluster import KMeans, DBSCAN, AgglomerativeClustering
from sklearn.metrics import silhouette_score
from sklearn.metrics import silhouette_score, davies_bouldin_score, calinski_harabasz
from scipy.cluster.hierarchy import linkage, dendrogram

x = pd.read_csv("UCI HAR Dataset/features.csv")
y = pd.read_csv("UCI HAR Dataset/activity_labels.csv")

print(x.head())
print(y.head())
print(x.shape)
print(y.shape)
print("Missing Values in X :", x.isnull().sum().sum())
print("Missing Values in Y :", y.isnull().sum().sum())
\end{lstlisting}

\section{Standardization and PCA}
\begin{lstlisting}[language=Python]
print(y[0].unique())

scaler = StandardScaler()
x_scaled = scaler.fit_transform(x)

pca = PCA(n_components=2)
x_pca = pca.fit_transform(x_scaled)

print(x_scaled.shape)
print(x_pca.shape)
\end{lstlisting}

\section{K-Means Clustering}
\begin{lstlisting}[language=Python]
kmeans = KMeans(n_clusters=6, random_state=42)
kmeans_labels = kmeans.fit_predict(x_scaled)

kmeans_sil = silhouette_score(x_scaled, kmeans_labels)
kmeans_db = davies_bouldin_score(x_scaled, kmeans_labels)
kmeans_ch = calinski_harabasz_score(x_scaled, kmeans_labels)

kmeans_ari = adjusted_rand_score(
    y.values.ravel(), kmeans_labels)

kmeans_nmi = normalized_mutual_info_score(
    y.values.ravel(), kmeans_labels)

print("Silhouette :", kmeans_sil)
print("Davies-Bouldin :", kmeans_db)
print("Calinski-Harabasz :", kmeans_ch)
print("ARI :", kmeans_ari)
print("NMI :", kmeans_nmi)
\end{lstlisting}

\subsection{K-Means Results}
\begin{verbatim}
Silhouette : 0.1086410631864108
Davies-Bouldin : 2.3579675526882222
Calinski-Harabasz : 1862.506263577485
ARI : 0.41971368603551157
NMI : 0.5589238295237167
\end{verbatim}

\subsection{Elbow Method}
\begin{lstlisting}[language=Python]
wcss = []
silhoue_score = []

for k in range(2,9,1):
    kmeans = KMeans(n_clusters=k, random_state=42)
    kmeans.fit(x_scaled)
    wcss.append(kmeans.inertia_)

    labels = kmeans.fit_predict(x_scaled)
    score = silhouette_score(x_scaled, labels)
    silhoue_score.append(score)

result_table = pd.DataFrame({
    "K": range(2,9),
    "WCSS": wcss,
    "Silhouette Score": silhoue_score
})

print(result_table)
\end{lstlisting}

\begin{table}[H]
\centering
\caption{K-Means Results}
\begin{tabular}{ccc}
\toprule
$K$ & WCSS & Silhouette Score\\
\midrule
2 & 2336223.79 & 0.39650\\
3 & 2081849.62 & 0.327408\\
4 & 1976475.19 & 0.161332\\
5 & 1875068.02 & 0.13064\\
6 & 1818790.06 & 0.10864\\
7 & 1776518.72 & 0.11142\\
8 & 1748685.34 & 0.08959\\
\bottomrule
\end{tabular}
\end{table}

\section{DBSCAN Clustering}
\begin{lstlisting}[language=Python]
dbscan = DBSCAN(eps=1.5, min_samples=10)
db_labels = dbscan.fit_predict(x_pca)

sb_cluster_list = set(db_labels)
print(sb_cluster_list)

noise = 0
for ele in sb_cluster_list:
    if ele == -1:
        noise = noise + 1

print(noise)
\end{lstlisting}

\begin{verbatim}
{0, 1, 2, 3, -1}
1
\end{verbatim}

\subsection{DBSCAN Evaluation}
\begin{verbatim}
Silhouette Score: 0.490960875197007
Davies-Bouldin Index: 1.377583115821548
Calinski-Harabasz Index: 6437.629875579848
ARI: 0.3266064560591039
NMI: 0.5230949748586399
\end{verbatim}

\section{Hierarchical Agglomerative Clustering}
\begin{lstlisting}[language=Python]
hac = AgglomerativeClustering(n_clusters=6, linkage='ward')
hac_labels = hac.fit_predict(x_scaled)

sample = x_scaled[:300]
Z = linkage(sample, method='ward')

plt.figure(figsize=(12,5))
dendrogram(Z)
plt.title("Hierarchical Clustering Dendrogram")
plt.xlabel("Samples")
plt.ylabel("Distance")
plt.show()
\end{lstlisting}

\subsection{HAC Evaluation}
\begin{verbatim}
Silhouette: 0.08339649239977742
Davies-Bouldin: 2.7857334406214562
Calinski-Harabasz: 1730.8358392108255
ARI: 0.2793864491481101
NMI: 0.4544001173140205
\end{verbatim}

\section{Algorithm Comparison}
\begin{table}[H]
\centering
\begin{tabular}{lccccc}
\toprule
Algorithm & Silhouette & Davies-Bouldin & Calinski-Harabasz & ARI & NMI\\
\midrule
K-Means & 0.108641 & 2.357968 & 1862.506264 & 0.419714 & 0.558924\\
DBSCAN & 0.490961 & 1.377583 & 6437.629876 & 0.326606 & 0.523095\\
Hierarchical & 0.083396 & 2.785733 & 1730.835840 & 0.279386 & 0.454400\\
\bottomrule
\end{tabular}
\caption{Comparison of Clustering Algorithms}
\end{table}

\section{Observations}
\begin{itemize}
\item K-Means was evaluated for $k=2$ to $8$. WCSS decreased as the number of clusters increased.
\item The highest Silhouette Score was obtained for $k=2$, with a score of 0.39650.
\item As $k$ increased, WCSS continued to decrease, while the Silhouette Score generally decreased.
\item PCA reduced the dimensionality of the dataset and made the data easier to visualize.
\end{itemize}

\section{Learning Outcomes}
\begin{itemize}
\item Understood the basic concept of unsupervised learning and clustering.
\item Learned to preprocess Human Activity Recognition data before clustering.
\item Understood how K-Means groups data points based on similarity.
\item Learned how the Elbow Method and WCSS can be used to study the suitable number of clusters.
\item Learned how the Silhouette Score evaluates cluster separation.
\item Understood the basic working of DBSCAN and how it can identify dense regions and noise points.
\end{itemize}

\end{document}
